\revision{Our first observation is that an extension to more than two populations should retain conjugacy. 
Analogous to our response to AE Comment 4, we can accomplish the conjugacy goal by including a series of proper beta distribution factors while also allowing a leading term  in front of two beta distributions:
\begin{equation}
    \nu(\mathrm{d}\btheta)= f(\btheta) \left[ \prod_{p=1}^{P} \theta_p^{\gamma_p-1}(1-\theta_p)^{c_p-1} \right] \mathrm{d}\btheta.
\end{equation}
}

\revision{We want to satisfy Desiderata (A) and (B) while also making sure $f(\btheta)$ does not factorize in to separate functions of the elements of $\btheta$.
To that end, we first recall the $P=2$ case in which $\sigma_1=\sigma_2=\sigma$. 
In that special two-population case, the leading term we chose becomes $(\theta_1+\theta_2)^{-\sigma-\gamma_1-\gamma_2}$. It is straightforward to extend this term to $P$ populations by choosing $(\sum_{p=1}^{P}\theta_p)^{-\sigma-\sum_{p=1}^P \gamma_p}$.
Importantly, this construction satisfies BNP desiderata (A) and (B); we observe that the same machinery we employ for the 2-dimensional proofs should also hold for the $P$-dimensional generalization.
In particular, we can bound $\max_p\{\theta_p\}\le \sum_{p=1}^{P}\theta_p\le P\max_p\{\theta_p\}$.}

\revision{The remaining issue then is how to allow potentially different power law growth rates in each population.
To see how to proceed, first recall $P=2$ but now where $\sigma_1$ need not equal $\sigma_2$. In this general case with $P=2$, our leading term $f(\btheta)$ was $(\theta_1+\theta_2^{\sigma_2/\sigma_1})^{\sigma_1}(\theta_1+\theta_2)^{-\gamma_1-\gamma_2}$. We can extend this leading term to a general number of populations $P$ by choosing $f(\btheta) = (\theta_1+\sum_{p=2}^P \theta_p^{\sigma_p/\sigma_1})^{\sigma_1}(\sum_{p=1}^P \theta_p)^{-\sum_{p=1}^P \gamma_p}$. We expect the resulting rate measure to retain all of: conjugacy, Desiderata (A) and (B), and the potential for separate power law growth rates across populations.}

